\documentclass{exam}

\newif\ifprintselected   % required by exam_config's \ansbox

%%%%%%%%%%%%%%%%%%%%%%%%%%%%%%%%%%%%%%%%%%%%%%%%%%%%%%
%%% SET THESE IF VARIABLES
%%%%%%%%%%%%%%%%%%%%%%%%%%%%%%%%%%%%%%%%%%%%%%%%%%%%%%
% v2 of the midterm packet, written for Fall 2027. One page on purpose: the
% v1 packet (ECE444_Project_Midterm.tex, 8 pages) prescribed every step, and
% real projects arrive with a goal and a deadline, not a procedure. The
% decisions v1 made for the cadets (the match bar, which planes, the f2
% window, the balun) are theirs to make and defend here.
\ifdefined\thisTermIn
	\def\thisterm{\thisTermIn}
\else
	\def\thisterm{Fall 2027}
\fi
\def\thisclass{ECE 444}

\def\fOne{915}

% Due date and weights. Both are quoted from the syllabus -- change there first.
\def\duedate{Lesson 20, 2359}
\def\projweight{30\%}
\def\midweight{40\%}

%%%%%%%%%%%%%%%%%%%%%%%%%%%%%%%%%%%%%%%%%%%%%%%%%%%%%%
%%% END SET THESE IF VARIABLES
%%%%%%%%%%%%%%%%%%%%%%%%%%%%%%%%%%%%%%%%%%%%%%%%%%%%%%
\input{myPackages_gr}
\usepackage[hidelinks, linkcolor=red]{hyperref}
\input{myShortcuts}
% ECE 444 document fonts.
%
% Body text is Barlow (SIL OFL; the faces and the license are in
% latex/fonts/Barlow/). Barlow is not in TeX Live, so it is loaded as an
% OpenType face through fontspec -- which means the document must be built
% with lualatex, not pdflatex. build_practice.sh and build_lab.sh both do.
%
% Compiling with pdflatex anyway is not an error: the guard below falls back
% to the LaTeX default so the source still compiles under any engine and can
% be syntax-checked. The resulting PDF just is not the one we ship.
%
% Paths are relative to latex/, which is where both build scripts run
% (each does cd "$(dirname "$0")"). Loading by path rather than by installed
% family name means no system font install and no fontconfig, so a fresh
% clone builds a byte-comparable PDF anywhere.
%
% MATH IS DELIBERATELY LEFT IN COMPUTER MODERN. Barlow has no math companion,
% and CM math under sans body text is already the house pairing: the decks set
% words in Source Sans Pro and let MathJax render symbols in its CM-derived
% face. Moving math to a sans as well would mean unicode-math plus a sans math
% font (Fira Math being the plausible one) and re-checking every equation in
% the course for symbol coverage. That is a separate decision, not a side
% effect of this one.

\usepackage{iftex}
\ifLuaTeX
	\usepackage{fontspec}
	\defaultfontfeatures{Ligatures=TeX}
	\setmainfont{Barlow}[
		Path           = fonts/Barlow/,
		Extension      = .ttf,
		UprightFont    = *-Regular,
		ItalicFont     = *-Italic,
		BoldFont       = *-Bold,
		BoldItalicFont = *-BoldItalic,
	]
	\setsansfont{Barlow}[
		Path           = fonts/Barlow/,
		Extension      = .ttf,
		UprightFont    = *-Regular,
		ItalicFont     = *-Italic,
		BoldFont       = *-SemiBold,
		BoldItalicFont = *-BoldItalic,
	]
	% microtype earns its place here specifically because of Barlow. Its
	% glyphs are narrower than Computer Modern's, so justified prose sets
	% tighter and TeX has to squeeze interword space to fit; two paragraphs
	% in the L02/L03 sets went overfull on the first Barlow build. Protrusion
	% and font expansion absorb that -- both are lualatex features, so this
	% belongs inside the guard with everything else.
	\usepackage{microtype}

	% Mono stays Latin Modern: Barlow ships no monospaced face, and verbatim
	% and code need one that actually is fixed-pitch.
\fi

\setlength{\parindent}{0pt}
\setlength{\parskip}{5pt}

\input{exam_config}
\noprintanswers
\printselectedfalse

% Compact run-in headings: \section's spacing alone pushed v2 onto a second page.
\newcommand{\head}[1]{\par\vspace{6pt}{\large\bfseries #1}\par\vspace{-2pt}}

\firstpageheader{}{}{}
\runningheader{\thisclass \hspace{1pt} -- \thisterm }{Midterm Project}{\thepage\hspace{1pt} of\hspace{1pt} \numpages}
\firstpagefooter{}{}{}
\runningfooter{}{}{}
\runningheadrule

\begin{document}
\pagestyle{headandfoot}

\begin{center}
USAF ACADEMY \\
DEPARTMENT OF ELECTRICAL AND COMPUTER ENGINEERING \\
MIDTERM PROJECT (assigned Lesson 9) - Dipole Design, Matching, and Measurement \\
\thisterm \\[2mm]
\end{center}

\textbf{Due \duedate.} Individual. Worth \projweight\ of the course grade and
\midweight\ of your midterm grade. There is no resubmission.

\head{The Assignment}

Build a dipole, make it work at the frequency it was built for, then make the
same antenna work at a frequency it was not built for. Measure what each
approach does to the impedance and to the pattern, and write it up.

\begin{enumerate}[nosep]
	\item Design and build a dipole for \fOne\ MHz.
	\item Tune it by trimming.
	\item Characterize it with S-parameters and input impedance.
	\item Measure its pattern in the chamber, co-polarized and cross-polarized.
	\item Pick a different frequency and measure the antenna's impedance there.
	      Leave the antenna as it is.
	\item Design and build an L-network that matches it at that frequency.
	\item Measure the matched antenna's S-parameters and its pattern, co-polarized
	      and cross-polarized.
	\item Write a report detailing your design, your process, and your analysis.
\end{enumerate}

\head{What Is Up to You}

This description is short on purpose. Real projects arrive with a goal and a
deadline, not a procedure, and most of the engineering is in the decisions the
goal leaves open. What counts as ``tuned''? Which second frequency, and why that
one? Which cuts, over what span, and how do you know your setup is not part of
the measurement? Make those calls, write down why you made them, and defend
them in the report. Lessons 4, 7, and 9--11 have the tools; deciding which ones
apply is part of the work.

Where a decision runs into the limits of our equipment (the chamber, the
analyzers, the parts bin), ask me early. A question asked at Lesson 12 is free;
the same question asked the night before it is due is not.

\head{The Report}

One typed technical report, submitted as a PDF, with your measured data files
attached. There is no template and no page limit. Write it for an engineer who
was not in the lab with you: what you designed and why, how you built and
measured it, what you found, and what you make of it.

Compare what you measured against what you expected. You are not required to
make the numbers agree, but you are required to know why they do not. A
disagreement you have explained is a result; a disagreement you have not
mentioned is a missing result.

\head{Grading}

\begin{center}
\begin{tabular}{|p{4.9in}|c|}
\hline
\textbf{What is assessed} & \textbf{Points} \\
\hline
\textbf{Analysis.} Measured against expected at both frequencies, with every
disagreement explained or owned as unexplained. & 35 \\
\hline
\textbf{Measurements.} Impedance, S-parameters, and co- and cross-polarized
patterns, before and after the match, taken so that a reader can trust them. & 25 \\
\hline
\textbf{Design decisions.} The dipole, the trim, the second frequency, and the
network, each with its reasoning stated. & 20 \\
\hline
\textbf{Report quality.} Organized, typed, every figure captioned with labeled
axes and units, data files included. & 20 \\
\hline
\end{tabular}
\end{center}

\textbf{Collaboration.} You may share chamber sessions and help each other take
data. The design, the data reduction, the analysis, and the report are
individual. Generative AI use is Level 3; document it as the syllabus requires.

\end{document}
